\section{Logical operations and fresh quantification}
\label{sec:predicate-calculus}

Finite support is a property of a predicate, not a condition for forming a
proposition in Lean. This separation allows ordinary logical reasoning and
quotient elimination to coexist with operations whose nominal interpretation
requires a support certificate. The distinction is particularly useful for
quantification: quantifying over a whole acted carrier and fixing one of its
elements have different support requirements.

\subsection{Supported predicates and their semantic views}

The Lean carrier \texttt{SupportedPred A X} stores an ordinary predicate
$P:X\to\mathrm{Prop}$ and a proof that some finite bound supports it.
It stores no chosen bound. Application returns a proposition, and pointwise
logical equivalence gives equality of predicate values by propositional
extensionality. Proof irrelevance ensures that changing the certificate,
including replacing a sufficient bound by a larger one, changes no value.
The atom and carrier universes are independent; the record stays in the
universe of $X$.

Named conversions identify this carrier with supported maps into discrete
truth values and with supported subsets of $X$. The full function-object view
is \texttt{predicateObject}, and the satisfying-set view sends $P$ to
$\{x\mid P(x)\}$. Inverse precomposition defines the predicate action;
its laws follow by injective transfer from the function-object action.
Both supported equivalences preserve the action, and the set view intertwines
inverse precomposition with direct image. In particular,
\[
 (\pi\act P)(\pi\act x)\ \longleftrightarrow\ P(x).
\]
Every view preserves and reflects each sufficient finite support bound.
Consequently the supported predicate carrier is nominal even when $X$ is
merely a permutation set. None of these constructions needs infinite atoms
or decidable atom equality.

Over infinite atoms these support correspondences also identify least
supports. For the full function object and the bare satisfying set, the
equations use the individual support certificates induced by $P$; they do
not assert that every predicate or subset is finitely supported. The maps
remain explicit conversions, leaving one automatic coercion to ordinary
$X\to\mathrm{Prop}$ and preserving the existing actions on bare functions
and sets.

\subsection{Boolean structure and the scope of quantification}

Let $X$ carry an action of $\Perm(A)$, without assuming every element of
$X$ finitely supported. Write $\mathcal P_{\mathrm{fs}}(X)$ for the
finitely supported subsets under direct image, equivalently the predicates
under inverse precomposition in \eqref{eq:predicate-action}.
Complement preserves and reflects every supporting bound, using classical
double negation for reflection. Each binary logical connective preserves a
common bound for its operands; intersection and union in particular are
supported by the union of separate bounds. Every constant truth value has
empty support. Thus these subsets form a Boolean subalgebra of the ordinary
powerset. Implication is complement followed by union; equivalence is the
conjunction of the two implications. These bounds are sufficient, not generally
least: intersecting any predicate with the empty set erases its dependence.

This formulation separates a nominal closure argument from ordinary order
theory. Mathlib's \texttt{BooleanSubalgebra} supplies the inherited Boolean
algebra once closure is established. The equivalence \texttt{subsetEquiv}
transfers that algebra to \texttt{SupportedPred} using
\texttt{Equiv.booleanAlgebra}. Its order and operations compute as ordinary
logic, while the supported-subset action is the restriction of the existing
set image action. No independent proof of the Boolean algebra axioms is needed.
Nor does this construction imply
closure under every externally indexed union.

The Lean operation \texttt{biimp} is derived from the two implications, and
the application and coerced-function equations expose these operations as
ordinary propositions. The action preserves every Boolean operation and
preserves and reflects the inclusion order. Over infinite atoms, complement
has exactly the original least support; the least support of a binary
operation is contained in the union of the operand supports.

For precomposition $P\circ f$, the ordinary map certificate for $f$ and the
predicate certificate for $P$ give their union as a sufficient bound.
The bundled operations \texttt{precomp} and \texttt{precompMap} express this
for an ordinary certified map and a supported map respectively. Fixing a
parameter $y$ of a jointly supported relation $R(y,x)$ likewise gives the
union of a bound for $R$ and a bound for this particular $y$. The operation
\texttt{section} accepts that individual certificate, even when the carrier
of $y$ contains unsupported elements. Simultaneously renaming the relation
and parameter renames its section. These constructions reuse the supported
function calculus; they introduce no new currying requirement.

\begin{proposition}[Quantification over an acted carrier]
\label{prop:predicate-quantification}
Let $X,Y$ be permutation sets. If a finite bound $S$ supports a relation
$R\subseteq X\times Y$ under the joint action, then the same bound supports
both predicates
\[
  x\longmapsto (\forall y\in Y)\,R(x,y),
  \qquad
  x\longmapsto (\exists y\in Y)\,R(x,y).
\]
Neither carrier need be nominal or nonempty, and the atom type need not be
infinite.
\end{proposition}

\begin{proof}
For a permutation $\pi$ fixing $S$, joint support gives
$R(\pi\act x,\pi\act y)\leftrightarrow R(x,y)$.
Reindex either quantifier by the bijection $y\mapsto\pi\act y$.
The argument also applies to an empty carrier.
\end{proof}

Restricted quantification incorporates the domain condition into the relation:
use $D(x,y)\Rightarrow R(x,y)$ for a universal quantifier and
$D(x,y)\wedge R(x,y)$ for an existential one. Only support of this combined
body is required; separate bounds for $D$ and $R$ are sufficient and combine
by finite union. Another sufficient condition for an external
family $(P_i)_{i\in I}$ is a single finite bound supporting every $P_i$;
no action or inhabitance on the index is needed in that case. In Lean the
index may be any \texttt{Sort}, in a universe independent of the atoms and
the acted carriers. This uniform condition
differs from joint support, which may rename the index, and from the weaker
assertion that each section has some finite support of its own.

In Lean, \texttt{SupportsPred.all} and \texttt{SupportsPred.ex} handle jointly
acted relations; \texttt{SupportsPred.iAll} and \texttt{SupportsPred.iEx} handle
uniform external families. Common-bound and existential finite-support laws
need no decidable atom equality premise. Displayed unions of bounds retain
that premise; existential proofs choose equality decisions locally.

The bundled counterparts \texttt{SupportedPred.all} and
\texttt{SupportedPred.ex} have precisely these ordinary quantified bodies.
They commute with renaming and retain every supplied bound of the relation;
over infinite atoms their least supports are therefore contained in the
relation's least support. For an empty quantified carrier the universal
result is the top predicate and the existential result is the bottom one.

For an infinite, coinfinite atom set $U$, every singleton $\{a\}$ with
$a\in U$ is supported, but their union $U$ is unsupported. There is not even
a least supported upper bound in the inclusion order: any supported superset
$V$ of $U$ is cofinite, and for $b\in V\setminus U$ the smaller set
$V\setminus\{b\}$ is still a supported upper bound. Such a $b$ exists because
$U$ is coinfinite. By contrast, a finitely supported
\emph{collection} of supported subsets has a supported union and intersection:
quantify over $\mathcal P_{\mathrm{fs}}(X)$ with the collection's membership
predicate as a guard, using the joint action on membership and evaluation.
Explicitly, for a supported collection $C$ these operations are
\[
 U_C(x)\ \longleftrightarrow\ (\exists P)\,(C(P)\wedge P(x)),
 \qquad
 I_C(x)\ \longleftrightarrow\ (\forall P)\,(C(P)\Rightarrow P(x)).
\]
If $\pi$ fixes a bound of $C$, then $C(\pi\act P)\leftrightarrow C(P)$,
while evaluation satisfies $(\pi\act P)(\pi\act x)\leftrightarrow P(x)$.
Thus both guarded relations have the same bound as $C$, and
Proposition~\ref{prop:predicate-quantification} gives that bound for $U_C$
and $I_C$. Reindexing the predicates also proves
$U_{\pi\act C}=\pi\act U_C$ and $I_{\pi\act C}=\pi\act I_C$.
The Lean operations \texttt{collectionUnion} and \texttt{collectionInter}
use this construction. Their least-support inclusions over infinite atoms
follow directly, without requiring $X$ nominal.
For nominal $X$, these are the supported-family operations in Pitts' nominal powerset
\cite[Section~2.5 and Example~11.1]{pitts2013}.

\subsection{Ordinary descent and reflection of predicate support}
\label{sec:predicate-descent}

Let $\sim$ be any equivalence relation on $X$, and write $q:X\to X/{\sim}$.
A raw predicate descends precisely when
\[
  x\sim y\quad\Longrightarrow\quad P(x)\leftrightarrow P(y).
\]
Consequently there is an ordinary correspondence
\[
 (X/{\sim}\to\mathrm{Prop})
 \simeq
 \{P:X\to\mathrm{Prop}\mid
       \forall x,y,\ x\sim y\Rightarrow(P(x)\leftrightarrow P(y))\}.
\]
One direction is precomposition with $q$. In Lean, propositional
extensionality turns compatibility into equality of truth values, permitting
\texttt{Quotient.lift}; its representative equation proves one inverse law,
and quotient induction proves the other. Mathlib's \texttt{Setoid.liftEquiv}
already expresses this universal property for ordinary functions, so the
predicate correspondence is its specialization through propositional
extensionality. No permutation action or finite-support certificate is involved.

\begin{proposition}[Support reflection by predicate pullback]
\label{prop:predicate-pullback}
Suppose $q:X\to Y$ is an equivariant surjection between permutation sets.
For every ordinary predicate $P:Y\to\mathrm{Prop}$ and finite atom bound $S$,
\[
  S\text{ supports }P\circ q
  \quad\Longleftrightarrow\quad
  S\text{ supports }P.
\]
Neither carrier is required to be nominal.
\end{proposition}

\begin{proof}
Preservation follows by substituting
$q(\pi\act x)=\pi\act q(x)$ in the support equation for $P$.
For reflection, write an arbitrary $y$ as $q(x)$ and make the same
substitution in the equation for $P\circ q$.
\end{proof}

The reflection mechanism only needs a scalar operation, its preservation by
$q$, and surjectivity; inversion and atom infinitude play no role. Without
surjectivity, the pullback cannot detect predicate behavior away from the image.
For a nonempty example take the diagonal $q:\mathbb N\to\mathbb N^2$ under
finite permutations of $\mathbb N$, and
$P(x,y)\equiv(x=0\wedge x\ne y)$. Its pullback is constantly false, but
swapping $0$ and $1$ changes its value at $(0,2)$.

The function-object version is equally general about bounds: if $G$ is a
division monoid and $q:X\to Y$ is an equivariant surjection, precomposition
embeds the full conjugation object $Y\to Z$ into $X\to Z$ equivariantly.
Compatibility with conjugation uses the action equation for $q$ at $g^{-1}$,
without cancellation. Injective action transfer therefore reflects every
sufficient set of an acted bound carrier, including infinite sets. This does
not extend least finite support minimality to arbitrary infinite bounds.

For an invariant equivalence relation, the canonical quotient action makes
the projection equivariant. Ordinary descent therefore preserves and reflects
each sufficient support bound. On supported predicates it also preserves least
support over infinite atoms. This differs from the support inclusion and fiber
intersection for quotient \emph{objects} in Section~\ref{sec:equivariant-quotients}:
it makes no assertion that a representative attains the support of its class.
An alpha-incompatible predicate, such as one that distinguishes representative
binder labels identified by the relation, fails the ordinary compatibility
condition before support is considered. For dependent data-valued motives,
transport coherence is a separate requirement; propositional extensionality
does not discharge that requirement for arbitrary fibers.

The Lean namespace \texttt{PredicateDescent} provides this correspondence
as \texttt{ordinaryEquiv}. The theorem \texttt{descend\_mk} computes on
representatives; both round trips are equalities of ordinary functions.
All support and renaming
statements select \texttt{QuotientAction.mulAction} explicitly. An unrelated
action on the same quotient type is not covered. Unsupported predicates can
still descend: among atoms $\mathbb N\times\mathrm{Bool}$, the true-flag
predicate is infinite and coinfinite, yet it descends through a quotient
forgetting an extra discrete label. Ordinary quotient induction can then
extract a true-flag representative from a proof of that predicate.

\subsection{The supported quotient correspondence}

Let $\sim$ be invariant under the selected action on $X$ and give
$X/{\sim}$ its canonical quotient action. Compatibility is stable under
inverse precomposition: if $x\sim y$, then
$\pi^{-1}\act x\sim\pi^{-1}\act y$, so a compatible $P$ gives
$(\pi\act P)(x)\leftrightarrow(\pi\act P)(y)$.
The compatible supported predicates therefore inherit a restricted action.
The ordinary correspondence restricts to an equivariant equivalence
\[
 \mathcal P_{\mathrm{fs}}(X/{\sim})
 \simeq
 \{P\in\mathcal P_{\mathrm{fs}}(X)\mid P\text{ is compatible with }\sim\}.
\]
Both directions retain exactly the sufficient bounds, by
Proposition~\ref{prop:predicate-pullback}. The restricted carrier has the
same bounds as its underlying predicates and is nominal even when $X$ is
not nominal. Neither this construction nor its representative equations
requires infinite atoms.

In Lean, \texttt{CompatiblePred} is this subtype of \texttt{SupportedPred}.
The explicit constructor \texttt{compatibleAction} restricts the existing
action using Mathlib's \texttt{SubMulAction}. Nominality is supplied by the
explicit \texttt{compatibleNominal} certificate.
The descent operation is \texttt{descendSupported}; the equivalence is
\texttt{supportedEquiv}. Representative computation remains ordinary application:
\[
 \operatorname{descend}(P)([x])\ \longleftrightarrow\ P(x).
\]
Their inverse equations and action compatibility use the selected quotient
and restricted actions. Top, bottom, complement, intersection, union,
implication and logical equivalence are preserved, as follows by computation
on representatives and quotient induction. These are laws of the existing
Boolean operations on predicate values, with compatibility of the operands.

Over infinite atoms, equality of all sufficient finite bounds gives equality
of least supports. Thus a quotient predicate and its corresponding compatible
predicate have the same freshness relation to any individually supported
context $c$: disjointness with $\operatorname{supp}(c)$ is identical. The carrier containing
$c$ need not be nominal.

For example, quotient $A\times\Discrete_A(\mathrm{Bool})$ by equality of the first
coordinate. For fixed $a\in A$, the predicate $P(b,\ell)\equiv b=a$ is
compatible and supported by $\{a\}$. Every supplied finite
$S\supseteq\{a\}$ remains a bound after descent, including strictly larger
bounds. The descended predicate holds at both $[(a,\mathsf{true})]$ and
$[(a,\mathsf{false})]$ and, over infinite atoms, has least support $\{a\}$.
The extra labels affect neither predicate truth nor support.

\subsection{Cofinite truth and the Some/Any principle}
\label{sec:some-any}

For any ordinary atom predicate $P$, define
\[
  \mathsf N a.\,P(a)
  \quad\Longleftrightarrow\quad
  \{a\in A\mid\neg P(a)\}\text{ is finite}.
\]
This is eventual truth in Mathlib's cofinite filter, and is Pitts'
freshness quantifier \cite[Definition~3.8]{pitts2013}. Its definition needs
neither support nor infinitude. Infinitude ensures that the filter is proper:
an eventual property then has a witness. On a finite carrier every property
is cofinite, including the constantly false predicate.

The Lean operation is \texttt{Freshly}, with binder notation available in the
opt-in scope \texttt{NominalPackage.FreshQuantifier}. Its finite-exception
form says that there is a finite $S$ such that every $a\notin S$ satisfies
$P$. No action is used in this equivalence. Cofinite truth is preserved and
reflected by any equivalence of input types, even in different universes.

\begin{proposition}[Some/Any for a supplied bound]
\label{prop:some-any-bound}
Assume $A$ is infinite and $S$ is any finite support of $P:A\to\mathrm{Prop}$
for the canonical atom action. Then
\[
  \mathsf N a.\,P(a)
  \quad\Longleftrightarrow\quad
  \exists a\notin S,\ P(a)
  \quad\Longleftrightarrow\quad
  \forall a\notin S,\ P(a).
\]
In particular, if $a\notin S$, the fresh-quantified proposition is equivalent
to $P(a)$.
\end{proposition}

\begin{proof}
For $a,b\notin S$, the transposition $(a\ b)$ fixes $S$, so support
gives $P(a)\leftrightarrow P(b)$. Thus one witness outside $S$ gives
every such witness. The universal statement implies cofinite truth because
all exceptions lie in $S$. Conversely, a cofinite truth set meets the
complement of $S$, since their intersection is cofinite in an infinite set.
\end{proof}

This is the sufficient-bound form of Pitts' Lemma~3.7; his contextual
Some/Any theorem follows by taking a supported section of a jointly
equivariant relation \cite[Theorem~3.9]{pitts2013}.
More generally, a relation bound and a bound for an individual parameter
combine by union. No nominality of the whole parameter carrier is needed.
Enlarging the bound with any finite avoidance set leaves the principle valid;
the proof never computes an exact least support. Fresh quantification also
preserves the support of a jointly supported relation, by reindexing the
cofinite filter along the atom permutation. The ordinary theorem
\texttt{SupportsPred.fresh} retains the supplied bound for
$x\mapsto\mathsf N a.\,R(x,a)$ without atom infinitude or nominality of $X$.

The bundled operation \texttt{SupportedPred.fresh} has this ordinary body
and commutes with renaming. Every sufficient bound of $R$ supports its fresh
projection; over infinite atoms this gives
$\operatorname{supp}(R.\mathsf{fresh})\subseteq\operatorname{supp}(R)$.
For an invariant relation and an individually supported $x$, the context form is
\[
 \mathsf N a.\,R(x,a)
 \quad\longleftrightarrow\quad
 \exists a\mathrel{\#}x,\ R(x,a)
 \quad\longleftrightarrow\quad
 \forall a\mathrel{\#}x,\ R(x,a).
\]
An arbitrary finite exclusion can be imposed in both quantified forms.
Here the freshness premise uses the individual certificate's least support:
Lean writes \texttt{hx.Fresh a}, without requiring the carrier of $x$ nominal.
For example, the identity is supported in the full conjugation carrier of
functions between permutation groups even though that carrier contains
unsupported constant functions. It is therefore a valid parameter in this
principle. Atom predicates similarly admit evaluation using freshness for
their individually certified function objects, or for their bundled supported
values. These are consequences of the existing support-disjointness relation.

Under the canonical atom action, a predicate is finitely supported precisely
when its truth set is finite or cofinite. This classification needs no
infinitude: constancy outside a supplied bound gives one of the two cases,
while a finite truth set supports its membership predicate and complementation
preserves support. Consequently every supported atom predicate is eventually
decided: $\mathsf N a.\,P(a)$ or $\mathsf N a.\,\neg P(a)$.

The logical laws have different hypotheses. Monotonicity, conjunction and
modus ponens are filter laws for arbitrary predicates. Self-dual negation
requires both infinitude and eventual decision. Decision of either operand
suffices for disjunction, and decision of the antecedent suffices for implication:
\[
 \mathsf N a.\,(P(a)\Rightarrow Q(a))
 \quad\longleftrightarrow\quad
 ((\mathsf N a.\,P(a))\Rightarrow(\mathsf N a.\,Q(a))).
\]
These latter two laws need no infinitude; the other operand may be arbitrary.
For the analogous equivalence law, decision of both operands suffices.
Alternatively an already cofinite operand suffices with an arbitrary other
operand. A merely supported operand alone does not suffice. These refinements separate the filter
argument from the nominal hypotheses sufficient to use it; Pitts states
the Boolean laws for supported operands \cite[Proposition~3.10]{pitts2013}.

The valid one-way quantifier laws are
\[
 (\mathsf N a.\,\forall i,\,P_i(a))\Rightarrow
       (\forall i,\,\mathsf N a.\,P_i(a)),\qquad
 (\exists i,\,\mathsf N a.\,P_i(a))\Rightarrow
       (\mathsf N a.\,\exists i,\,P_i(a)).
\]
Finite indices give the converse universal law. For an arbitrary external
index, a common finite support bound for all $P_i$ also gives that converse,
without infinitude. The common bound gives the converse existential law if
the atoms are infinite, or alternatively if the index is nonempty. In the
infinite case, choose one atom outside the common bound and use Some/Any on
every section; in the finite case all predicates are cofinite. The empty-index
case explains the additional existential premise: over finite atoms cofinite
falsehood holds, while an existential over an empty index fails. Lean permits
these external indices in any \texttt{Sort}, with no action on them.

Arbitrary quantifier interchange remains invalid even for jointly equivariant
relations. The statement $\mathsf N a.\,\exists x,\ a=x$ holds, whereas
$\exists x,\mathsf N a.\,a=x$ fails. Dually,
$\forall x,\mathsf N a.\,a\ne x$ holds but
$\mathsf N a.\,\forall x,\ a\ne x$ fails.
Finite nonempty indices do not repair existential interchange: among atoms
$\mathbb N\times\mathrm{Bool}$, the predicate $P(n,b)\equiv b=\mathsf{true}$
and its complement have cofinite union, while neither is cofinite. Both are
unsupported by the classification above. The same example refutes unrestricted
self-dual negation, disjunction and implication. Pairing $P$ with the constantly
false, supported predicate refutes the equivalence law with only one supported
operand: both individual cofinite assertions are false, but their pointwise
equivalence is not cofinite.

Nor can atom Some/Any be transferred to every nominal carrier. On
$\Discrete_{\mathbb N}(\mathrm{Bool})\times\mathbb N$, the predicate selecting
the true discrete label has empty support and a witness, but has infinitely
many exceptions. The semantic role of the canonical atom action is essential.

Finally, finite fresh-name existence does not give a finitely supported global
selector. Suppose $f:\mathcal P_{\mathrm{fin}}(A)\to A$ satisfies $f(S)\notin S$
and has a finite map-support bound $T$. Since $T$ supports itself, it supports
$a=f(T)$. Set $b=f(T\cup\{a\})$. Then $a,b\notin T$ and $a\ne b$, so their
transposition fixes $T$ but moves $a$, contradicting its support. This argument
needs no infinitude assumption: the proposed selector supplies the spare atom.
Ordinary classical choice of a fresh atom on an infinite type remains legal;
the obstruction concerns its finite support.
